\documentclass[10pt]{article}
\usepackage[T1]{fontenc}
\usepackage[utf8]{inputenc}
\usepackage[french]{babel}
\usepackage{amsmath,amssymb,mathrsfs,mathpazo}
\usepackage{geometry}
\geometry{a4paper,margin=2cm}
\usepackage{graphicx,array,tabularx,tabularray,colortbl,multirow,multicol,booktabs,enumitem,pifont,ulem,fancybox,diagbox,cancel}
\usepackage{pstricks,pstricks-add,pst-plot,pst-tree,pst-node,pst-text,pst-func,pst-math,pst-bspline,pst-3dplot,pst-eucl,pst-solides3d}
\let\pscalculate\undefined
\usepackage{pst-bezier}
\usepackage{tikz,pgfplots,tkz-tab}
\usetikzlibrary{arrows,arrows.meta,patterns,calc,positioning,trees,decorations.markings,shapes.geometric,shapes.misc,shapes.arrows,intersections,decorations.pathmorphing,decorations.pathreplacing,backgrounds,shadows,babel}
\pgfplotsset{compat=1.18}
\usepgfplotslibrary{fillbetween}
\usepackage[np]{numprint}
\usepackage{listings,makecell,pas-tableur,twoopt}
\let\barre\undefined
\usepackage[pstricks]{tablvar}
\definecolor{darkgreen}{RGB}{0,100,0}
\providecommand{\red}{\color{red}}
\providecommand{\blue}{\color{blue}}
\providecommand{\green}{\color{green}}
\providecommand{\black}{\color{black}}
\providecommand{\white}{\color{white}}
\providecommand{\gray}{\color{gray}}
\providecommand{\cyan}{\color{cyan}}
\providecommand{\magenta}{\color{magenta}}
\providecommand{\yellow}{\color{yellow}}
\providecommand{\violet}{\color{violet}}
\providecommand{\orange}{\color{orange}}
\providecommand{\euro}{\texteuro}
\providecommand{\EUR}{\texteuro}
\providecommand{\ieme}{\textsuperscript{e}}
\providecommand{\degres}{\ensuremath{^\circ}}
\providecommand{\degree}{\ensuremath{^\circ}}
\providecommand{\vv}[1]{\overrightarrow{#1}}
\providecommand{\Coord}[1]{#1}
\providecommand{\decosix}{}
\providecommand{\starredbullet}{\ensuremath{\bullet}}
\pagestyle{empty}
\setlength{\parindent}{0pt}
\usepackage{adjustbox}\DeclareUnicodeCharacter{2236}{:}\DeclareUnicodeCharacter{2190}{\ensuremath{\leftarrow}}
\providecommand{\R}{\mathbb{R}}
\providecommand{\N}{\mathbb{N}}
\providecommand{\Z}{\mathbb{Z}}
\providecommand{\Q}{\mathbb{Q}}
\providecommand{\D}{\mathbb{D}}
\providecommand{\C}{\mathbb{C}}
\providecommand{\e}{{\mathrm{e}}}
\providecommand{\vect}[1]{\overrightarrow{#1}}
\providecommand{\vectt}[1]{\overrightarrow{\mathrm{#1}}}
\providecommand{\vvt}[1]{\overrightarrow{\mathrm{#1}}}
\providecommand{\np}[2][]{{#2~#1}}
\providecommand{\vv}[1]{{\overrightarrow{#1}}}
\renewcommand{\d}{{\mathrm{d}}}
\providecommand{\degres}{{^\circ}}
\providecommand{\degre}{{^\circ}}
\providecommand{\blue}{\color{blue}}
\providecommand{\red}{\color{red}}
\providecommand{\enskip}{\quad}
\providecommand{\llbracket}{\lbrack\!\lbrack}
\providecommand{\rrbracket}{\rbrack\!\rbrack}
\providecommand{\euro}{{\text{€}}}
\providecommand{\no}{{\text{n°}}}
\providecommand{\interval}[2]{{[#1;#2]}}
\providecommand{\intervFO}[2]{{[#1;#2[}}
\providecommand{\dotfill}{\ldots}

\providecommand{\ds}{\displaystyle}
\providecommand{\textding}[1]{#1}
\providecommand{\Oij}{\(\mathrm{O};\vec{\imath},\vec{\jmath}\)}
\providecommand{\Oijk}{\(\mathrm{O};\vec{\imath},\vec{\jmath},\vec{k}\)}
\providecommand{\Ouv}{\(\mathrm{O};\vec{u},\vec{v}\)}
\providecommand{\intervFF}[2]{\ensuremath{[#1;#2]}}
\renewcommand{\euro}{\ifmmode\text{\texteuro}\else\texteuro\fi}
\DeclareUnicodeCharacter{00B0}{\ifmmode^{\circ}\else\textdegree\fi}
\DeclareUnicodeCharacter{20AC}{\ifmmode\text{\texteuro}\else\texteuro\fi}
% Énoncés : Éducation nationale. Transcription et corrigés : APMEP, Denis Vergès et contributeurs. Découpage et classement : CoopMaths, Sébastien Lozano, Éric Elter et contributeurs. Adaptation de la présentation : Maths Explorer.
% CC BY-SA : https://coopmaths.fr/www/about/
% https://coopmaths.fr/alea/static/bac/2024/tex/bac_2024_11_sujet2_amsud_3_cor.tex
\begin{document}
\noindent\textit{APMEP / CoopMaths. CC BY-SA. Présentation adaptée par Maths Explorer.}\par\medskip
\medskip

\textbf{Partie 1}

\medskip

On considère la fonction $f$ définie sur l'ensemble des nombres réels $\R$ par :
\[f(x) = \left(x^2 - 4\right)\e^{-x}.\]



\begin{enumerate}
\item Limites :

$\bullet~$ : en $- \infty$ : de $\displaystyle\lim_{x \to - \infty}x^2 = + \infty$ d'où 
$\displaystyle\lim_{x \to - \infty}x^2 - 4 = + \infty$ et $\displaystyle\lim_{x \to - \infty}\e^{-x} =
 + \infty$ on en déduit par produit de limites que $\displaystyle\lim_{x \to - \infty} f(x) = + \infty$ ;

$\bullet~$ : en $+ \infty$ $f(x) = x^2\e^{-x} - 4\e^{-x}$ : on sait que 

$\displaystyle\lim_{x \to + \infty}4\e^{-x} = 0$ et que $\displaystyle\lim_{x \to + \infty}x^2\e^{-x} = \displaystyle\lim_{x \to + \infty}\dfrac{x^2}{\e^x} =  0$ (puissances comparées), donc :

$\displaystyle\lim_{x \to + \infty} f(x) = 0$.
\item 
$f$ est dérivable comme produit de deux fonctions dérivables sur $\R$. Sur cet intervalle :

$f'(x) = 2x\e^{-x} + \left(x^2 - 4\right) \times (- 1)\e^{-x} = \e^{-x}(2x - x^2 + 4) = \left(- x^2 + 2x + 4\right)\e^{-x}$.
\item 
On sait que quel que soit $x \in\R,\, \e^{-x} > 0$, donc le signe de $f'(x)$ est celui du trinôme $- x^2 + 2x + 4$.

On a $- x^2 + 2x + 4 = -\left(x^2 - 2x - 4\right) = -\left[(x - 1)^2 - 1 - 4\right] = -\left[(x - 1)^2 - 5\right] =$

$- x^2 + 2x + 4 = - \left(x - 1 + \sqrt 5\right)\left(x - 1 - \sqrt 5\right)$ : ce trinôme a deux racines $1 - \sqrt 5$ et $1 + \sqrt 5$.

On sait que le trinôme a le signe de $a = -1$, donc est négatif sauf sur l'intervalle 

$\left[1 - \sqrt 5~;~1 + \sqrt 5\right]$ où il est positif.

On a donc le tableau de variations :

\begin{center}
\psset{unit=1cm,arrowsize=2pt 3}
\begin{pspicture}(10,3)
\psframe(10,3)\psline(0,2)(10,2)\psline(0,2.5)(10,2.5)\psline(1,0)(1,3)
\uput[u](0.5,2.4){$x$} \uput[u](1.4,2.4){$-\infty$} \uput[u](4,2.4){\footnotesize$1 - \sqrt 5$} \uput[u](5.5,2.4){$0$} \uput[u](7,2.4){\footnotesize$1 + \sqrt 5$} \uput[u](9.6,2.4){$+ \infty$} 
\rput(0.5,2.25){$f'(x)$}\rput(0.5,1){$f$}
\psline{->}(1.5,1.5)(3.5,0.5)\psline{->}(4.5,0.5)(6.5,1.5)\psline{->}(7.5,1.5)(9.5,0.5)
\rput(2.5,2.25){$-$}\rput(5.5,2.25){$+$}\rput(8.5,2.25){$-$}
\rput(5.5,1){$-4$}\uput[d](1.4,2){$+ \infty$}\uput[u](4,0){$\approx -8,5$}\uput[d](7,2){$\approx 0,25$}
\uput[u](9.4,0){$0$}
\end{pspicture}
\end{center}

On a $f\left(1 - \sqrt 5\right) = \left(2 - 2\sqrt 5 \right)\e^{\sqrt 5 - 1} \approx - 8,5$ ;

$f\left(1 + \sqrt 5\right) = \left(2 + 2\sqrt 5 \right)\e^{\sqrt 5 -+ 1} \approx 0,25$ ;

$f(0) = - 4$
\end{enumerate}

\bigskip

\textbf{Partie 2}

\medskip

On considère la suite $(I_n)$ définie pour tout entier naturel $n$ par $I_n = \displaystyle\int_{-2}^0 x^n\e^{-x}\, \text{d}x$.

\medskip

\begin{enumerate}
\item 
$I_0 = \displaystyle\int_{-2}^0 \e^{-x}\, \text{d}x = \left[- \e^{-x}\right]_{-2}^0 = - \e^{0} - \e^{2} = \e^2 - 1$.
\item 
On calcule $I_n$ en faisant une intégration par partie : on a 

$\left\{\begin{array}{l c l l c l}
u(x) &=& \e^{-x}&u'(x) &=& - \e^{-x}\\
v'(x)&=&x^n&v(x) &=&\dfrac{x^{n+1}}{n+1}\\
\end{array}\right.$ on a 

$I_n = \left[\e^{-x} \times \dfrac{x^{n+1}}{n+1}\right]_{-2}^0
 - \displaystyle\int_{-2}^0 - \e^{-x} \times \dfrac{x^{n+1}}{n+1}\, \text{d}x = 0 - \dfrac{(-2)^{n+1}\e^{2}}{n + 1} + \dfrac{I_{n+1}}{n + 1} \iff $
 
 $(n + 1)I_n = -(- 2)^{n+1}\e^2 + I_{n+1} \iff I_{n+1} = (- 2)^{n+1} \text{e}^2 + (n + 1)I_n.$


\item
L'égalité précédente donne avec :

$\bullet~n = 0 :\, I_{1} = (- 2)^{1}\text{e}^2 + I_0 = - 2\e^{2} + \e^{2} - 1 = -\e^{2} - 1$,

$\bullet~n = 1 :\, I_{2} = (- 2)^{2} \text{e}^2 + I_1 = 4 \text{e}^2  - 2\e^{2} - 2 = 2\text{e}^2 - 2$.
\end{enumerate}

\medskip

\textbf{Partie 3}

\medskip

\begin{enumerate}
\item 
$f$ a le signe du trinôme $x^2 - 4$ car quel que soit $x \in \R$, \, $\e^{- x} > 0$.

Comme $x^2 - 4 = (x + 2)(x - 2)$ ce trinôme a le signe de $a = 1$ donc est positif, sauf sur l'intervalle borné par les deux racines soit sur l'intervalle $]-2~;~2[$ où $f'x) < 0$.

Donc sur l'intervalle $]-2~;~2[, \, f(x) < 0$ et sur $]- \infty~;~-2[ \cup ]2~;~+ \infty[, \, f(x) > 0$.
\item 




$\mathcal{C}_f$ coupe l'axe des abscisses  au point d'abscisse $- 2$ t on vient de voir que sur l'intervalle $]-2~,~0[, \, f(x) < 0$, donc l'aire $S$ du domaine $D$ est égale à l'opposé de l'intégrale de la fonction $d$ de $x = - 2$ à $x = 0$.

$\text{aire}(D) = - \displaystyle\int_{-2}^0 f(x)\:\text{d}x = - \displaystyle\int_{-2}^0 \left(x^2 - 4\right)\e^{-x}\:\text{d}x = - \displaystyle\int_{-2}^0 x^2\e^{-x}  + 4\displaystyle\int_{-2}^0 \e^{-x} = - I_2 + 4I_0 = 4\left(-\e^{2} - 1\right) - \left(2\e^2 - 2\right) = 2\e^2 - 2 \approx 12,78$ (ce que l'on peut conforter avec la figure).
\end{enumerate}

\begin{center}
\psset{unit=1cm}
\begin{pspicture*}(-3,-9)(5,3)
\psgrid[gridlabels=0,subgriddiv=1,gridwidth=0.2pt](-2,-9)(0,0)
\psline{->}(0,0)(1,0)\psline{->}(0,0)(0,1)\uput[d](1,0){\scriptsize $\vect{\imath}$}
\uput[l](0,1){\scriptsize $\vect{\jmath}$}
\psaxes[Dx=10,Dy=20]{->}(0,0)(-3,-9)(5,3)
\psplot[plotpoints=2000,linewidth=1.pt,linecolor=red]{-3}{6}{x dup mul 4 sub 2.71828 x exp div}
\pscustom[fillstyle=hlines,linecolor=red]{\psplot[plotpoints=2000,linewidth=1.2pt,linecolor=red]{-2}{0}{x 2 exp  4 sub 2.71828 x exp div}\psline[linewidth=0.2pt](0,-4)(0,0)(-2,0)}\uput[dl](0,0){\small O}
\uput[r](-2,2){\red $\mathcal{C}_f$}\rput(-1,-3){$D$}
\end{pspicture*}
\end{center}

\bigskip

\end{document}
